\section{Recomendaciones}

\begin{enumerate}
	\item Instalar un sensor de presión diferencial entre el laboratorio y la zona adyacente, con un set-point de control entre +12.5 Pa y +25 Pa, para garantizar y monitorear la presurización positiva de forma continua.
	\item Filtrar el aire de impulsión con filtros de eficiencia MERV 13-14 (decisión definitiva). El laboratorio es de análisis industrial y no de biocontención, por lo que no se requiere filtración HEPA; los escenarios HEPA evaluados quedan descartados.
	\item Instalar un damper de alivio calibrable como componente obligatorio del lazo de control de presión: a la densidad del sitio (0.88 kg/m$^3$) las rejillas a 3.0 m/s sostienen naturalmente solo 11 Pa, por debajo del mínimo admisible de 12.5 Pa, de modo que el set-point de +25 Pa solo se garantiza modulando la exfiltración con el damper y el sensor de presión diferencial \cite{dml_hd_inst_001,dml_listado_equipos_2026}.
	\item Especificar materiales resistentes al ambiente altamente corrosivo de la planta de hipoclorito de calcio: ventilador, rejillas, damper y soportes en PRFV o acero inoxidable 316, o con recubrimientos epóxicos, y motor TEFC anticorrosivo.
	\item Sellar adecuadamente puertas, ventanas, pasos de instalaciones y cualquier abertura no controlada del recinto, con burletes y umbrales herméticos, para reducir las fugas no controladas y facilitar el mantenimiento de la presión positiva.
	\item Ubicar las tres rejillas de exfiltración en paredes opuestas o perpendiculares a la impulsión, de modo que se favorezca un flujo cruzado que minimice zonas de estancamiento.
	\item Validar la presión diferencial real en campo tras el balanceo y ajuste del damper de alivio, comparándola con el campo de presión obtenido en el modelo CFD. Considerar un estudio de edad del aire o trazadores si se requiere demostrar la eficacia de mezcla en zonas con recirculación.
	\item Prever un margen de sobredimensionamiento del 10 \% al 20 \% entre el caudal de impulsión y cualquier caudal de extracción localizada (campanas, puntos de captación), de modo que la presión positiva general del recinto no se vea comprometida.
\end{enumerate}
